\thispagestyle{empty}
\begin{center}
\begin{LARGE}
\textbf{Resumen}
\end{LARGE}
\end{center}
\begin{quotation}

La identificaci�n de patrones de motricidad asociados al Trastorno del Espectro Autista (TEA) mediante la observaci�n de videos constituye un desaf�o que requiere herramientas complementarias para facilitar el an�lisis de las conductas observadas. Ante esta problem�tica, el objetivo del presente trabajo es desarrollar un sistema web asistido por inteligencia artificial para identificar patrones de motricidad relacionados con el TEA, espec�ficamente aleteo de manos, balanceo corporal y movimientos repetitivos de cabeza, como apoyo a la evaluaci�n profesional, sin sustituir el diagn�stico cl�nico.
La metodolog�a comprendi� la selecci�n y organizaci�n de videos, el etiquetado de las conductas observadas, la extracci�n de fotogramas y el preprocesamiento de im�genes mediante normalizaci�n y generaci�n de secuencias. Para el reconocimiento de los movimientos, se desarrollaron y compararon dos modelos de aprendizaje profundo: una red neuronal convolucional (CNN) propia combinada con una red de memoria a corto y largo plazo (LSTM), y una arquitectura basada en ResNet50 + LSTM. Posteriormente, el modelo con mejor desempe�o se integr� en un prototipo web que permite cargar videos, analizarlos y visualizar los patrones identificados junto con informaci�n complementaria.
La evaluaci�n experimental, realizada sobre un conjunto de prueba de 15 muestras, mostr� que ResNet50 + LSTM obtuvo una exactitud del 73,33 % y un F1-Score ponderado del 74,00 %, mientras que la CNN propia alcanz� una exactitud del 60,00 % y un F1-Score ponderado del 57,00 %. Estos resultados permitieron seleccionar ResNet50 + LSTM para su integraci�n en la plataforma. Sin embargo, el tama�o reducido del conjunto de prueba limita el alcance de las conclusiones y evidencia la necesidad de ampliar los datos y realizar evaluaciones adicionales. El sistema constituye una herramienta tecnol�gica de apoyo para el an�lisis de patrones de motricidad y una base para futuras investigaciones, sin emitir diagn�sticos cl�nicos ni reemplazar el criterio de los profesionales de la salud.
\vspace*{0.5cm}

\noindent {\bf Descriptores:} 1. $<$descriptor1$>$, 2. $<$descriptor2$>$, 3. $<$descriptor3$>$.

\end{quotation}
